%% kapa: exam=1 round=24 year=2013 period=1 form=A subject=economics subject_ko=경제원론 seq=38 number=78
%% source: Q-Net 감정평가사 기출문제 (cst00309) artlSeq=5073802 file=extracted/1008/5073802/2013년 제24회 감정평가사 1차시험 기출문제/1교시(A형).hwp
%% conversion: kordoc 4.24.51 -> tools/gen1.py (auto); review=proofread-sonnet; flags=none
\begin{question}{78}
甲은 乙에게 100만 원을 위탁하고, 乙은 자금을 운용하여 이익이 발생할 때에는 이익의 10\%를 운용수수료로 받고 손실이 발생할 때에는 운용수수료를 받지 않는 계약을 맺었다. 甲은 자금운용을 전적으로 乙에게 위임하며 乙은 자신의 이익만을 추구한다. 상황에 따른 각 투자안별 수익(원금 포함)이 다음과 같을 때 乙의 투자선택과 기대운용수수료는? (단, 甲과 乙은 위험중립적이다.) \par
{\small\centering
\begin{tabular}{|c|c|c|}
\hline
 & \makecell{상황 1} & \makecell{상황 2} \\
\hline
\makecell{확률} & \makecell{0.1} & \makecell{0.9} \\
\hline
\makecell{X 투자안} & \makecell{100만 원} & \makecell{110만 원} \\
\hline
\makecell{Y 투자안} & \makecell{400만 원} & \makecell{70만 원} \\
\hline
\end{tabular}\par}

\begin{choices}
  \item X 투자안을 선택하며, 기대운용수수료는 3만 원이다.
  \item X 투자안을 선택하며, 기대운용수수료는 6만 원이다.
  \item X 투자안을 선택하며, 기대운용수수료는 9만 원이다.
  \item Y 투자안을 선택하며, 기대운용수수료는 3만 원이다.
  \item Y 투자안을 선택하며, 기대운용수수료는 6만 원이다.
\end{choices}
\end{question}
\answer{4}
